\chapter{What Linear Algebra Is Really About}

In the Preface, we promised pictures before formulas -- so before we compute a determinant or
row-reduce a matrix, let us agree on what this whole subject is \emph{for}. Most people meet linear algebra as a pile of
procedures: solve this system, invert this matrix, find these eigenvalues. The procedures are real and we
will learn them carefully. But if that is all you take away, you will have memorized the grammar of a
language without ever hearing it spoken. In this book we start with the geometric
\emph{pictures}, so that each algebraic procedure has a clear visual meaning.

So here is the one-sentence version, which we will spend the rest of the chapter unpacking:

\begin{focusbox}
Linear algebra is the mathematics of \textbf{vectors}, the \textbf{linear combinations} you can
build from them, and the \textbf{linear transformations} that move whole spaces around while keeping
straight lines straight and the origin fixed. Almost everything else is a consequence of those three
ideas.
\end{focusbox}

If that sentence is unfamiliar, think of a \textbf{vector} as a displacement or an arrow
pointing from here to there; think of a \textbf{linear combination} as a recipe that scales a few arrows
and stitches them tip to tail; and think of a \textbf{linear transformation} as a rule that reshapes the
entire plane or space while keeping grid lines straight and evenly spaced. By the end of this chapter, these
three words will each carry a sharp, instinctive mental image, and you will understand why the adjective
\emph{linear} is the most useful kind of structure a piece of mathematics can have.

\section{What Makes Linearity Special}

Mathematics is full of functions. Some are wild: $\sin x$ wiggles forever, $e^x$ explodes,
$1/x$ blows up at the origin. Linear algebra deliberately restricts attention to well-behaved
operations -- and discovers that they govern a wide range of mathematics and
practical applications.

What does it mean for an operation to be \emph{linear}? There are just
two fundamental moves you are allowed to make: you may \textbf{scale} (stretch or shrink a thing by a number) and
you may \textbf{add} (combine two things tip to tail). A linear operation is one that does not care
about the order in which you do these: scaling-then-transforming gives the same answer as
transforming-then-scaling, and the same for adding. Grid lines stay straight and evenly spaced.

\begin{intuitionbox}
Picture a function $f$ as a machine. It is \emph{linear} if it respects two promises:
\begin{itemize}
  \item \textbf{Scaling:} if you feed it twice as much, you get twice as much out. $f(2\vec v) = 2f(\vec v)$.
  \item \textbf{Adding:} if you feed it $\vec u$ plus $\vec v$, you get $f(\vec u)$ plus $f(\vec v)$.
\end{itemize}
A linear machine cannot tell the difference between ``combine first, then process'' and ``process
each piece, then combine.'' That single property is what makes linear systems analytically tractable.
\end{intuitionbox}

\begin{figure}[H]
\centering
\begin{tikzpicture}[scale=0.85,>=Stealth]
  % Panel 1: Original Cartesian grid
  \begin{scope}[shift={(0,0)}]
    \draw[step=0.6,gray!25,very thin] (-1.5,-1.5) grid (1.5,1.5);
    \draw[->,gray!60] (-1.6,0) -- (1.6,0) node[right]{\scriptsize $x$};
    \draw[->,gray!60] (0,-1.6) -- (0,1.6) node[above]{\scriptsize $y$};
    \draw[->,very thick,vecblue] (0,0) -- (0.6,0) node[below=1pt]{\scriptsize $\vec{e}_1$};
    \draw[->,very thick,vecred] (0,0) -- (0,0.6) node[pos=0.6,left=2pt]{\scriptsize $\vec{e}_2$};
    \filldraw[black] (0,0) circle (1.5pt);
    \node at (0,-2.1) {\small (a) Original grid};
  \end{scope}

  % Panel 2: Linear transformation (shear + stretch)
  \begin{scope}[shift={(5.2,0)}]
    \foreach \k in {-2,...,2} {
      \draw[gray!35,thin] ({0.6*\k - 1.2*0.3},{-1.2}) -- ({0.6*\k + 1.2*0.3},{1.2});
      \draw[gray!35,thin] ({-1.4 + 0.6*\k*0.3},{0.6*\k}) -- ({1.4 + 0.6*\k*0.3},{0.6*\k});
    }
    \draw[->,gray!60] (-1.6,0) -- (1.6,0);
    \draw[->,gray!60] (0,-1.6) -- (0,1.6);
    \draw[->,very thick,vecblue] (0,0) -- (0.6,0) node[below=1pt]{\scriptsize $\mat{A}\vec{e}_1$};
    \draw[->,very thick,vecred] (0,0) -- (0.3,0.6) node[above left=1pt,fill=white,inner sep=1pt]{\scriptsize $\mat{A}\vec{e}_2$};
    \filldraw[black] (0,0) circle (1.5pt);
    \node[targetgreen!60!black,align=center] at (0,-2.1) {\small (b) Linear: lines stay straight,\\\small parallels stay parallel, $\vec{0}$ fixed};
  \end{scope}

  % Panel 3: Nonlinear transformation (curved lines)
  \begin{scope}[shift={(10.6,0)}]
    \foreach \k in {-2,...,2} {
      \draw[gray!35,thin,domain=-1.2:1.2,samples=20] plot ({\x + 0.2*sin(120*\x)}, {0.6*\k + 0.15*cos(120*\x)});
      \draw[gray!35,thin,domain=-1.2:1.2,samples=20] plot ({0.6*\k + 0.18*sin(120*\x)}, {\x + 0.15*sin(90*\x)});
    }
    \draw[->,gray!60] (-1.6,0) -- (1.6,0);
    \draw[->,gray!60] (0,-1.6) -- (0,1.6);
    \filldraw[trapred] (0.2,0.25) circle (1.5pt) node[above right=1pt,fill=white,inner sep=1pt]{\scriptsize $f(\vec{0})\neq\vec{0}$};
    \node[trapred!80!black,align=center] at (0,-2.1) {\small (c) Nonlinear: lines curve,\\\small spacing warps, origin wanders};
  \end{scope}
\end{tikzpicture}
\caption{The visual geometry of linearity. (a) A regular square Cartesian grid. (b) A linear transformation moves the grid but keeps all grid lines completely straight, keeps parallel lines parallel, and keeps the origin firmly pinned in place. (c) A nonlinear map bends and curves the grid, destroying the predictable combination structure.}
\label{ch01:linearity-grid}
\end{figure}

Why should we care so much about such a restrictive family? There are three main reasons.

\textbf{First, linear things are predictable.} If I know what a linear machine does to a few basic
inputs, I know what it does to \emph{every} input, because every input is a combination of the basic
ones and the machine respects combinations. You will see this turn into the entire theory of
matrices: a matrix is simply a record of ``what the machine did to the basic inputs,'' and that
record is enough to reconstruct everything.

\textbf{Second, the world is full of linear things, or things we choose to approximate as linear.}
The current through a circuit is a linear function of the voltages. The position of a robot arm is
(locally) a linear function of its joint angles. The output of one layer of a neural network is a
linear function of its input, wrapped in a simple nonlinear squish. Even when nature is curved,
calculus tells us that \emph{up close} everything looks linear -- a derivative is exactly a ``best
linear approximation.'' In multivariable calculus, that local linear approximation is represented by a matrix of derivatives.

\textbf{Third, linear problems are the ones we can actually solve.} We have complete,
reliable theories for linear systems. No comparably general theory exists for arbitrary nonlinear systems. So an enormous amount of applied mathematics follows a standard strategy:
approximate a difficult nonlinear problem by a sequence of linear problems. Newton's
method, gradient descent, the simplex algorithm, and the optimization routines that train modern machine
learning models all rely on this local linearization.

\begin{recoverybox}
\textbf{Pause and test your intuition:} Is the ordinary line $f(x) = 2x + 1$ a linear function in the sense
of linear algebra?

In everyday speech and high-school algebra, we call anything with a straight graph ``linear.'' But test our
two promises:
\begin{itemize}
  \item \textbf{Test scaling:} double the input from $x=1$ to $x=2$. At $x=1$, the function gives $f(1) = 2(1)+1 = 3$.
  If the machine scaled properly, doubling the input would double the output to $6$. But at $x=2$, we get
  $f(2) = 2(2)+1 = 5 \neq 6$. The scaling promise is broken!
  \item \textbf{Test the origin:} feed in zero. We get $f(0) = 1 \neq 0$. But any genuinely linear map must keep
  the origin nailed down, because $f(0\cdot x) = 0\cdot f(x) = 0$.
\end{itemize}
So $f(x) = 2x+1$ is what mathematicians call an \emph{affine} function -- a linear operation ($2x$) plus an
offset shift ($+1$). In linear algebra, a function is linear only if it preserves scaling and addition, which
forces its graph to pass directly through the origin.
\end{recoverybox}

\begin{aestheticbox}
``Linear'' sounds like a restriction -- no curves allowed!
-- yet that restriction is exactly what makes the theory comprehensive. We can completely classify
linear maps: determine when solutions exist, compute inverses, and decompose transformations into
rotations and scalings. Few branches of mathematics are so thoroughly understood.
\end{aestheticbox}

\section{Vectors: Things You Can Scale and Add}

Everything starts with the vector. You may have met vectors in physics as ``arrows with magnitude and
direction,'' and that picture is excellent -- we will use it constantly. But the deeper definition is
not about arrows at all. A vector is simply \textbf{a thing you are allowed to scale and add}, where
scaling and adding obey the standard rules. The arrow is one example; there are many others, and
recognizing them is a key skill.

Let us anchor the arrow picture first, because pictures come first in this book. Take the plane $\R^2$.
A vector $\vec v = \colvec{3\\1}$ is an instruction: ``go $3$ to the right and $1$ up.'' We draw it as
an arrow from the origin to the point $(3,1)$, as in Figure~\ref{ch01:arrow}.

\begin{figure}[H]
\centering
\begin{tikzpicture}[scale=1.0,>=Stealth]
  \draw[help lines,gray!30] (-1,-1) grid (5,3);
  \draw[->,gray!60] (-1,0)--(5,0) node[right]{$x$};
  \draw[->,gray!60] (0,-1)--(0,3) node[above]{$y$};
  \draw[->,vecblue,very thick] (0,0)--(3,1) node[midway,above,sloped]{$\vec v$};
  \filldraw[vecblue] (3,1) circle (1.5pt) node[above right=1pt]{$(3,1)$};
\end{tikzpicture}
\caption{A vector in $\R^2$ as an arrow from the origin. The vector $\vec v=\colvec{3\\1}$ encodes the
instruction ``$3$ right, $1$ up.''}
\label{ch01:arrow}
\end{figure}

\textbf{Scaling.} Multiply a vector by a number (a \emph{scalar}) and you stretch the arrow. The number
$2$ doubles its length; $\tfrac12$ halves it; $-1$ flips it to point the opposite way. The direction is
preserved (or exactly reversed); only the length changes.
\[
2\,\colvec{3\\1} = \colvec{6\\2},\qquad
-\colvec{3\\1} = \colvec{-3\\-1},\qquad
\tfrac13\,\colvec{3\\1} = \colvec{1\\\tfrac13}.
\]

\textbf{Adding.} To add two vectors, lay them tip to tail. Walk along the first, then continue along
the second; the sum is the single arrow from where you started to where you ended.
\[
\colvec{3\\1} + \colvec{1\\2} = \colvec{4\\3}.
\]
Figure~\ref{ch01:addscale} shows both operations side by side.

\begin{figure}[H]
\centering
\begin{tikzpicture}[scale=0.9,>=Stealth]
  \begin{scope}
    \draw[help lines,gray!25] (-1,-1) grid (5,3);
    \draw[->,gray!60] (-1,0)--(5,0); \draw[->,gray!60] (0,-1)--(0,3);
    \draw[->,vecblue,very thick] (0,0)--(3,1) node[midway,below]{$\vec u$};
    \draw[->,vecred,very thick] (3,1)--(4,3) node[midway,above right=1pt]{$\vec w$};
    \draw[->,structurepurple,ultra thick] (0,0)--(4,3) node[midway,above,sloped]{$\vec u+\vec w$};
    \node at (2.4,-1.6) {\small tip-to-tail addition};
  \end{scope}
  \begin{scope}[xshift=8cm]
    \draw[help lines,gray!25] (-3,-2) grid (4,3);
    \draw[->,gray!60] (-3,0)--(4,0); \draw[->,gray!60] (0,-2)--(0,3);
    \draw[->,vecblue,very thick] (0,0)--(3,1) node[above right=1pt]{\scriptsize $\vec v$};
    \draw[->,structurepurple,very thick] (0,0)--(1.5,0.5) node[above left=-1pt]{\scriptsize $\tfrac12\vec v$};
    \draw[->,vecred,very thick] (0,0)--(-3,-1) node[below left=1pt]{\scriptsize $-\vec v$};
    \node at (0.5,-2.6) {\small scaling stretches/flips};
  \end{scope}
\end{tikzpicture}
\caption{Left: vector addition is tip-to-tail. Right: scaling stretches ($\tfrac12\vec v$) or flips
($-\vec v$) without changing the line the vector lives on.}
\label{ch01:addscale}
\end{figure}

That is the entire arithmetic of vectors: scale and add. Notice we never multiplied two vectors
together to get a third vector -- that is deliberately \emph{not} part of the basic structure (it
comes later, and in more than one flavor). The plainness is the point.

\begin{formalbox}
A \textbf{real vector space} is a set $V$ together with an addition $\vec u + \vec v$ and a scalar
multiplication $c\,\vec v$ (with $c\in\R$) such that addition is commutative and associative, there is
a zero vector $\vec 0$ with $\vec v + \vec 0 = \vec v$, every $\vec v$ has a negative $-\vec v$, and
scaling distributes over both kinds of sum: $c(\vec u+\vec v)=c\vec u + c\vec v$ and
$(c+d)\vec v = c\vec v + d\vec v$, with $1\,\vec v = \vec v$. The familiar example is
$\R^n$, the set of columns of $n$ real numbers, but the definition deliberately mentions no arrows --
so that other things can qualify.
\end{formalbox}

\subsection{Vectors that are not arrows}

The abstract definition pays off immediately, because many different mathematical objects satisfy
the vector space axioms. Each example below is a collection of objects you can scale and add,
obeying the same rules, so the entire machinery of linear algebra applies directly to it.

\begin{itemize}
  \item \textbf{Lists of numbers.} A point in $\R^{784}$ is a $28\times 28$ grayscale image flattened
  into one long column. Scaling brightens it; adding overlays two images. Classifiers routinely
  treat images as vectors.
  \item \textbf{Functions.} You can add two functions ($(f+g)(x)=f(x)+g(x)$) and scale one
  ($(3f)(x)=3f(x)$). So functions form a vector space. This is the doorway to Fourier series and to
  all of functional analysis.
  \item \textbf{Polynomials.} The polynomials of degree at most $2$ form a $3$-dimensional vector
  space; $1$, $x$, and $x^2$ play the role of the basic directions.
  \item \textbf{Solutions of a linear differential equation.} If $y_1$ and $y_2$ both solve
  $y'' + y = 0$, so does $3y_1 - 2y_2$. The solution set is a vector space, which is exactly why such
  equations are so tractable.
\end{itemize}

\begin{connectionbox}
That functions are vectors is one of the great unifying ideas in mathematics. Once you accept
it, a Fourier series -- an infinite sum that breaks any periodic waveform into sinusoidal
harmonics -- is no longer an ad hoc trick; it is simply writing a function as a \emph{linear
combination} of the basic ``direction'' functions $\sin nx$ and $\cos nx$, exactly as you write
$\colvec{4\\3} = 4\colvec{1\\0} + 3\colvec{0\\1}$ in the plane. The infinite sum requires analysis to
make rigorous, but the underlying linear-algebraic idea is identical. Calculus, differential
equations, signal processing, and quantum mechanics all rest on this single observation, and we
will return to it once the finite-dimensional picture is solid.
\end{connectionbox}

\section{Linear Combinations and Span: Reachable Territory}

Now combine the two operations. A \textbf{linear combination} of vectors $\vec v_1,\dots,\vec v_k$ is
any vector you can build by scaling each one and adding the results:
\[
c_1\vec v_1 + c_2 \vec v_2 + \dots + c_k \vec v_k,
\qquad c_1,\dots,c_k \in \R.
\]
The numbers $c_i$ are the ``amounts'' of each ingredient. This is \emph{the} central construction of
linear algebra; nearly every later concept in the subject builds directly on linear combinations.

The natural follow-up question is: starting from a given set of vectors, scaling and adding as much as
you like, \emph{where can you get to}? The set of all reachable points is called the \textbf{span}.

\begin{intuitionbox}
Think of your vectors as the only moves available in a game. A linear combination is one sequence of
moves; the span is the entire region of the board you can ever reach. Span is the answer to
``what's possible from here?''
\end{intuitionbox}

The geometry is easiest to see in $\R^2$, so start there.

\textbf{One nonzero vector spans a line.} Take $\vec v = \colvec{2\\1}$. Every multiple $c\vec v$ lands
somewhere on the line through the origin in the direction of $\vec v$. You can reach any point on that
line and no point off it. Its span is a one-dimensional subspace: a line through the origin.

\textbf{Two vectors usually span the whole plane.} Add $\vec w = \colvec{-1\\1}$. Now you have two
independent directions. By choosing how much of each you take, you can reach \emph{any} point in
$\R^2$. For instance, to reach $\colvec{1\\4}$ we solve $c\colvec{2\\1}+d\colvec{-1\\1}=\colvec{1\\4}$,
which gives $c=\tfrac53,\ d=\tfrac73$. The span is all of $\R^2$.

\textbf{But two vectors that point the same way span only a line.} If $\vec w$ had been
$\colvec{4\\2}=2\vec v$, it would add nothing new -- it lies on the line $\vec v$ already reaches. Two
vectors span a plane only when they point in genuinely different directions. This is your first
glimpse of \emph{linear independence}, the idea that some vectors are redundant.
Figure~\ref{ch01:span} contrasts the two cases.

\begin{figure}[H]
\centering
\begin{tikzpicture}[scale=0.85,>=Stealth]
  \begin{scope}
    \draw[help lines,gray!20] (-3,-3) grid (3,3);
    \draw[->,gray!50] (-3,0)--(3,0); \draw[->,gray!50] (0,-3)--(0,3);
    \draw[structurepurple,thick,dashed] (-3,-1.5)--(3,1.5);
    \draw[->,vecblue,very thick] (0,0)--(2,1) node[above left=1pt,vecblue]{$\vec v$};
    \node at (0,-3.7) {\small span of one vector: a line};
  \end{scope}
  \begin{scope}[xshift=8cm]
    \fill[structurepurple!12] (-3,-3) rectangle (3,3);
    \draw[help lines,gray!25] (-3,-3) grid (3,3);
    \draw[->,gray!50] (-3,0)--(3,0); \draw[->,gray!50] (0,-3)--(0,3);
    \draw[->,vecblue,very thick] (0,0)--(2,1) node[above left=1pt,vecblue]{$\vec v$};
    \draw[->,vecred,very thick] (0,0)--(-1,1) node[above left=1pt,vecred]{$\vec w$};
    \node at (0,-3.7) {\small span of two independent vectors: all of $\R^2$};
  \end{scope}
\end{tikzpicture}
\caption{Span is reachable territory. One vector reaches a line; two independent vectors reach the
whole plane; a vector that merely repeats a direction reaches nothing new.}
\label{ch01:span}
\end{figure}

The same story scales up. In $\R^3$, one vector spans a line, two independent vectors span a plane
through the origin, and three independent vectors span all of space. These lines and planes through
the origin are our first examples of \textbf{subspaces} -- flat sheets, passing through $\vec 0$, that
are closed under scaling and adding. Subspaces are where vectors live, and recognizing them will let
us answer ``how many solutions does this system have?'' with a glance instead of a calculation.
Figure~\ref{ch01:subspace3d} shows the two proper subspaces of $\R^3$: a line and a plane, both
threaded through the origin.

\begin{figure}[H]
\centering
\begin{tikzpicture}[scale=1.0,>=Stealth]
  \begin{scope}
    % a plane through the origin in R^3
    \fill[structurepurple!18] (-1.6,-1.0) -- (1.6,-0.2) -- (1.6,1.4) -- (-1.6,0.6) -- cycle;
    \draw[structurepurple!70] (-1.6,-1.0) -- (1.6,-0.2) -- (1.6,1.4) -- (-1.6,0.6) -- cycle;
    \draw[->,gray!55] (0,0)--(1.8,0.55) node[right]{\scriptsize $y$};
    \draw[->,gray!55] (0,0)--(0,1.7) node[above]{\scriptsize $z$};
    \draw[->,gray!55] (0,0)--(-1.0,-0.7) node[below left,fill=white,inner sep=0.5pt]{\scriptsize $x$};
    \draw[->,vecblue,very thick] (0,0)--(1.1,0.2) node[below right=-1pt]{\scriptsize $\vec u$};
    \draw[->,vecred,very thick] (0,0)--(-0.7,0.5) node[above left=-1pt]{\scriptsize $\vec v$};
    \filldraw[black] (0,0) circle (1.3pt);
    \node at (0,-1.9){\small two independent vectors span a plane};
  \end{scope}
  \begin{scope}[xshift=6.5cm]
    \draw[->,gray!55] (0,0)--(1.8,0.55) node[right]{\scriptsize $y$};
    \draw[->,gray!55] (0,0)--(0,1.7) node[above]{\scriptsize $z$};
    \draw[->,gray!55] (0,0)--(-1.0,-0.7) node[below left]{\scriptsize $x$};
    \draw[structurepurple,very thick,dashed] (-1.4,-0.9)--(1.4,0.9);
    \draw[->,vecblue,very thick] (0,0)--(1.0,0.64) node[above left=-1pt]{\scriptsize $\vec v$};
    \filldraw[black] (0,0) circle (1.3pt);
    \node at (0,-1.9){\small one vector spans a line};
  \end{scope}
\end{tikzpicture}
\caption{Subspaces of $\R^3$ pass through the origin. Two independent vectors sweep out a plane
(left); a single nonzero vector sweeps out a line (right). Add a third independent vector and the
plane fills out to all of $\R^3$.}
\label{ch01:subspace3d}
\end{figure}

\begin{walkbox}
Take a walk and turn this over. In $\R^3$, you hold three vectors. Sometimes they span all of space,
sometimes only a plane, sometimes only a line. \emph{What geometric thing has to be true for three
vectors to fail to fill space?} Try to describe it without any formula -- in terms of arrows lying
flat. You are inventing the idea of linear dependence on your own, and it will feel obvious once you
say it out loud.
\end{walkbox}

\section{The Two Views of a Matrix: Rows and Columns}

We have been talking about vectors floating freely. Linear algebra gets its real power when we package
vectors into a grid -- a \textbf{matrix} -- and ask what a system of equations is really saying. There
are two complementary ways to read any linear system, and switching between them fluently is one of
the most valuable skills this book can give you. Strang famously builds his whole course around getting
you to see both. One small system will serve as our laboratory:
\[
\begin{aligned}
2x - y &= 0,\\
x + y &= 3.
\end{aligned}
\qquad\Longleftrightarrow\qquad
\underbrace{\begin{bmatrix} 2 & -1 \\ 1 & 1 \end{bmatrix}}_{\mat A}
\underbrace{\colvec{x\\y}}_{\vec x}
=
\underbrace{\colvec{0\\3}}_{\vec b}.
\]

\subsection{The row picture: intersecting lines}

Read the system one equation at a time, and each equation is a constraint. The first equation,
$2x - y = 0$, is satisfied by all points $(x,y)$ on a particular line. The second, $x + y = 3$, picks
out another line. A solution to the \emph{system} must satisfy both, so it must lie on \emph{both}
lines at once -- it is their intersection, drawn in Figure~\ref{ch01:rowpic}.

\begin{figure}[H]
\centering
\begin{tikzpicture}[scale=0.85,>=Stealth]
  \draw[help lines,gray!25] (-1,-1) grid (4,4);
  \draw[->,gray!55] (-1,0)--(4,0) node[right]{$x$};
  \draw[->,gray!55] (0,-1)--(0,4) node[above]{$y$};
  \draw[vecblue,very thick,domain=-0.6:2.1] plot (\x,{2*\x}) node[right]{$2x-y=0$};
  \draw[vecred,very thick,domain=-0.6:3.6] plot (\x,{3-\x}) node[right]{$x+y=3$};
  \filldraw[structurepurple] (1,2) circle (2.2pt);
  \node[structurepurple, font=\small] at (1.45, 2.2) {$(1,2)$};
\end{tikzpicture}
\caption{The row picture. Each equation is a line; the solution is where the lines cross, at
$(x,y)=(1,2)$.}
\label{ch01:rowpic}
\end{figure}

The two lines cross at $(1,2)$, so $x=1,\ y=2$. Geometrically a solution is an \emph{intersection of
hyperplanes} -- lines in $\R^2$, planes in $\R^3$, and higher-dimensional flats beyond. This view makes
the possible outcomes intuitive: two lines usually meet in one point, but parallel lines never meet
(no solution) and identical lines meet everywhere (infinitely many). Every strange thing a linear
system can do is, in the row picture, just lines arranging themselves.

\subsection{The column picture: mixing ingredients}

Now read the \emph{same} matrix equation a completely different way -- down the columns. Multiplying
$\mat A \vec x$ is, by definition, taking $x$ of the first column plus $y$ of the second:
\[
\colvec{x\\y}\ \longmapsto\
x\colvec{2\\1} + y\colvec{-1\\1}.
\]
So the equation $\mat A\vec x = \vec b$ is asking: \emph{what combination of the two column vectors
produces $\vec b = \colvec{0\\3}$?} This is the span question from the previous section, in a new
guise. The unknowns $x,y$ are no longer coordinates of a mystery point; they are the \emph{amounts}
of each column we must mix. Figure~\ref{ch01:colpic} shows the winning recipe.

\begin{figure}[H]
\centering
\begin{tikzpicture}[scale=0.85,>=Stealth]
  \draw[help lines,gray!25] (-2,-1) grid (3,4);
  \draw[->,gray!55] (-2,0)--(3,0) node[right]{};
  \draw[->,gray!55] (0,-1)--(0,4) node[above]{};
  % column 1 = (2,1), used once
  \draw[->,vecblue,very thick] (0,0)--(2,1) node[midway,above left=1pt,vecblue]{$\colvec{2\\1}$};
  % column 2 = (-1,1), used twice, tip to tail
  \draw[->,vecred,very thick] (2,1)--(0,3) node[midway,above right]{$2\colvec{-1\\1}$};
  \draw[->,structurepurple,ultra thick] (0,0)--(0,3) node[midway,left]{$\vec b=\colvec{0\\3}$};
\end{tikzpicture}
\caption{The column picture. We seek the recipe $x\colvec{2\\1}+y\colvec{-1\\1}=\colvec{0\\3}$. One
scoop of the blue column plus two scoops of the red column lands exactly on $\vec b$, so $x=1,\ y=2$.}
\label{ch01:colpic}
\end{figure}

And indeed $1\cdot\colvec{2\\1} + 2\cdot\colvec{-1\\1} = \colvec{0\\3}$. Same answer, $x=1,\ y=2$,
arrived at by a totally different mental route. The row picture intersected two lines; the column
picture mixed two ingredient-vectors. Both are correct, and a fluent reader flips between them
depending on which makes the current question easy.

\begin{focusbox}
Two ways to read $\mat A\vec x = \vec b$, and you want both fluent:
\begin{itemize}
  \item \textbf{Row view:} each row is an equation, a hyperplane; the solution is where all the
  hyperplanes intersect. Best for: ``does a solution exist, and is it unique?''
  \item \textbf{Column view:} $\mat A\vec x$ is a linear combination of $\mat A$'s columns; solving
  asks which recipe of columns builds $\vec b$. Best for: ``what can this matrix produce?'' -- the
  span of the columns, which we will name the \emph{column space}.
\end{itemize}
\end{focusbox}

\begin{trapbox}
The single most common early mistake is to think $\mat A\vec x$ means ``multiply each number of
$\mat A$ by each number of $\vec x$ somehow.'' It does not. \textbf{$\mat A\vec x$ is a linear
combination of the columns of $\mat A$, with the entries of $\vec x$ as the weights.} Burn the column
picture into your mind now and matrix multiplication will never again feel like an arbitrary rule. We
will derive the row-by-column ``dot product'' shortcut later, but it is a \emph{consequence} of the
column picture, not the meaning of it.
\end{trapbox}

\begin{labbox}
Do this one by hand before reading on. Take $\mat A=\begin{bmatrix}1&2\\3&1\end{bmatrix}$ and
$\vec b=\colvec{5\\5}$.
\begin{enumerate}
  \item \textbf{Row picture:} write the two equations, draw the two lines, and read off where they
  cross.
  \item \textbf{Column picture:} draw the columns $\colvec{1\\3}$ and $\colvec{2\\1}$ as arrows, and
  find by eye roughly how many scoops of each you need to reach $\colvec{5\\5}$.
  \item Check that both pictures give the same $(x,y)$. (Solving $x+2y=5$ and $3x+y=5$ gives
  $x=1,\ y=2$; one scoop of the first column and two of the second indeed land on $\colvec{5\\5}$.)
\end{enumerate}
Seeing one answer emerge from two unrelated-looking pictures is the moment linear algebra clicks.
\end{labbox}

\section{Linear Transformations: Matrices That Move Space}

So far a matrix has been a static array of coefficients. The deepest shift in viewpoint -- the one this
whole book is organized around -- is to see a matrix as a \textbf{verb}, not a noun. A matrix
\emph{does something}: it takes in a vector and spits out a vector. It is a function from space to
space. And not just any function -- a \textbf{linear transformation}, obeying the two promises from the
very start of this chapter.

\begin{intuitionbox}
Stop thinking of $\mat A\vec x$ as ``a calculation'' and start thinking of it as ``$\mat A$ \emph{acts
on} $\vec x$.'' The matrix $\mat A$ grabs the whole plane and moves it -- stretching, rotating,
shearing, flipping -- and $\mat A\vec x$ is wherever the point $\vec x$ ends up. Linearity means the
grid lines stay straight and evenly spaced, and the origin stays nailed down. Squares may become
parallelograms, but they never become curved blobs.
\end{intuitionbox}

The cleanest way to see this: a linear transformation is completely determined by what it does
to the basic direction vectors $\vec e_1 = \colvec{1\\0}$ and $\vec e_2 = \colvec{0\\1}$. Why?
Because every vector is a combination of them, $\colvec{x\\y} = x\vec e_1 + y\vec e_2$, and a linear
map respects combinations:
\[
\mat A\colvec{x\\y} = \mat A(x\vec e_1 + y\vec e_2) = x\,(\mat A\vec e_1) + y\,(\mat A\vec e_2).
\]
So if you tell me where the two basic arrows land, I can find where \emph{any} arrow lands. And those
two landing spots are precisely the \emph{columns} of the matrix. That is what a matrix really is: a
compact record of where the basic directions go.

\subsection{A gallery of transformations}

Watch four matrices act on the unit square (the grid spanned by $\vec e_1$ and $\vec e_2$),
collected in Figure~\ref{ch01:gallery}. Each picture shows the square before and after the matrix
grabs it.

\begin{figure}[H]
\centering
\begin{tikzpicture}[scale=0.88,>=Stealth]
  % rotation
  \begin{scope}
    \draw[help lines,gray!20] (-1.2,-1.2) grid (1.4,1.4);
    \draw[->,gray!55] (-1.2,0)--(1.4,0); \draw[->,gray!55] (0,-1.2)--(0,1.4);
    \fill[vecblue!15] (0,0)--(1,0)--(1,1)--(0,1)--cycle;
    \fill[vecred!18] (0,0)--(0.71,0.71)--(0,1.42)--(-0.71,0.71)--cycle;
    \draw[->,vecred,thick] (0,0)--(0.71,0.71) node[right=1pt]{\scriptsize $\mat{A}\vec{e}_1$};
    \draw[->,vecred,thick] (0,0)--(-0.71,0.71) node[above left=1pt]{\scriptsize $\mat{A}\vec{e}_2$};
    \node at (0,-1.8){\small rotation $45^\circ$};
  \end{scope}
  % scaling
  \begin{scope}[xshift=4.2cm]
    \draw[help lines,gray!20] (-0.4,-0.4) grid (2.2,1.6);
    \draw[->,gray!55] (-0.4,0)--(2.2,0); \draw[->,gray!55] (0,-0.4)--(0,1.6);
    \fill[vecblue!15] (0,0)--(1,0)--(1,1)--(0,1)--cycle;
    \fill[vecred!18] (0,0)--(2,0)--(2,0.5)--(0,0.5)--cycle;
    \draw[->,vecred,thick] (0,0)--(2,0) node[below=1pt]{\scriptsize $\mat{A}\vec{e}_1$};
    \draw[->,vecred,thick] (0,0)--(0,0.5) node[above left=1pt]{\scriptsize $\mat{A}\vec{e}_2$};
    \node at (0.9,-1.2){\small scale $x\!\times\!2$, $y\!\times\!\tfrac12$};
  \end{scope}
  % shear
  \begin{scope}[xshift=8.4cm]
    \draw[help lines,gray!20] (-0.4,-0.4) grid (2.4,1.6);
    \draw[->,gray!55] (-0.4,0)--(2.4,0); \draw[->,gray!55] (0,-0.4)--(0,1.6);
    \fill[vecblue!15] (0,0)--(1,0)--(1,1)--(0,1)--cycle;
    \fill[vecred!18] (0,0)--(1,0)--(2,1)--(1,1)--cycle;
    \draw[->,vecred,thick] (0,0)--(1,0) node[below=1pt]{\scriptsize $\mat{A}\vec{e}_1$};
    \draw[->,vecred,thick] (0,0)--(1,1) node[above right=1pt]{\scriptsize $\mat{A}\vec{e}_2$};
    \node at (0.9,-1.2){\small shear};
  \end{scope}
  % reflection
  \begin{scope}[xshift=12.6cm]
    \draw[help lines,gray!20] (-1.2,-0.4) grid (1.4,1.6);
    \draw[->,gray!55] (-1.2,0)--(1.4,0); \draw[->,gray!55] (0,-0.4)--(0,1.6);
    \fill[vecblue!15] (0,0)--(1,0)--(1,1)--(0,1)--cycle;
    \fill[vecred!18] (0,0)--(-1,0)--(-1,1)--(0,1)--cycle;
    \draw[->,vecred,thick] (0,0)--(-1,0) node[below=1pt]{\scriptsize $\mat{A}\vec{e}_1$};
    \draw[->,vecred,thick] (0,0)--(0,1) node[above left=1pt]{\scriptsize $\mat{A}\vec{e}_2$};
    \node at (0,-1.2){\small reflect across $y$-axis};
  \end{scope}
\end{tikzpicture}
\caption{Four linear transformations acting on the unit square (blue) and producing the image (red).
In every case straight lines stay straight, parallel lines stay parallel, and the origin stays put --
the signature of linearity. The red arrows are the new homes of $\vec e_1$ and $\vec e_2$, i.e.\ the
columns of the matrix.}
\label{ch01:gallery}
\end{figure}

The matrices that produce these are
\[
\underbrace{\begin{bmatrix} \tfrac{\sqrt2}{2} & -\tfrac{\sqrt2}{2}\\[2pt] \tfrac{\sqrt2}{2} & \tfrac{\sqrt2}{2}\end{bmatrix}}_{\text{rotation}},\quad
\underbrace{\begin{bmatrix} 2 & 0\\ 0 & \tfrac12\end{bmatrix}}_{\text{scaling}},\quad
\underbrace{\begin{bmatrix} 1 & 1\\ 0 & 1\end{bmatrix}}_{\text{shear}},\quad
\underbrace{\begin{bmatrix} -1 & 0\\ 0 & 1\end{bmatrix}}_{\text{reflection}}.
\]
Read each one column by column and confirm the picture. In the shear, for instance, $\vec e_1=\colvec{1\\0}$
stays put (first column $\colvec{1\\0}$) while $\vec e_2=\colvec{0\\1}$ slides over to $\colvec{1\\1}$
(second column) -- which is exactly the slanting you see.

\begin{example}
Trace one all the way through. Apply the shear $\mat S = \begin{bmatrix}1&1\\0&1\end{bmatrix}$ to
$\vec x=\colvec{2\\3}$ two ways.

\emph{Column (combination) view:}
\[
\mat S\colvec{2\\3} = 2\colvec{1\\0} + 3\colvec{1\\1} = \colvec{2\\0}+\colvec{3\\3}=\colvec{5\\3}.
\]
\emph{Row (dot-product shortcut) view:} top row $\inner{(1,1)}{(2,3)} = 1\cdot2+1\cdot3 = 5$; bottom
row $\inner{(0,1)}{(2,3)} = 0\cdot2+1\cdot3 = 3$. Same answer, $\colvec{5\\3}$. The point $(2,3)$ slid
horizontally by an amount equal to its height -- which is precisely what a shear does.
\end{example}

\begin{strangbox}
Strang opens his course by hammering on the row picture and the column picture of $\mat A\vec x=\vec b$,
and he is right to: that double vision is the foundation. This book agrees and goes one step further
\emph{before} you meet him. We frame the matrix as a \emph{transformation} -- a verb acting on space --
from the very first chapter, because that view is what later makes eigenvalues (``directions the
transformation merely stretches''), the SVD (``every transformation is rotate-stretch-rotate''), and
determinants (``the factor by which areas grow'') feel like one connected story instead of four
disconnected tricks. When you reach Strang's lectures, his row and column pictures will already be two
familiar faces of the single idea you started with here. And \emph{after} Strang, when you return for
the SVD and data chapters, this transformation lens is the thread that ties it all together.
\end{strangbox}

\subsection{Composition is matrix multiplication}

One more central result: if you apply transformation $\mat B$ and then transformation
$\mat A$, the combined effect ``$\mat A$ after $\mat B$'' is itself a linear transformation -- and the
matrix that represents it is exactly the product $\mat A\mat B$. This explains why matrix
multiplication is defined as it is: it is designed so that composing two linear transformations
corresponds to multiplying their matrices. Function composition \emph{is} matrix multiplication. We
will derive the coordinate formulas from this geometric principle in Chapter~5.

Figure~\ref{ch01:composition} shows the geometry: applying $\mat B$ reshapes the space once, and
applying $\mat A$ reshapes the result. The composite map that takes the starting grid straight to the
final destination in one leap is $\mat A\mat B$.

\begin{figure}[H]
\centering
\begin{tikzpicture}[scale=0.9,>=Stealth]
  % Stage 1: Unit Square
  \begin{scope}[shift={(0,0)}]
    \draw[help lines,gray!25] (-0.3,-0.3) grid (1.6,1.6);
    \draw[->,gray!60] (-0.4,0)--(1.7,0);
    \draw[->,gray!60] (0,-0.4)--(0,1.7);
    \fill[vecblue!20] (0,0)--(1,0)--(1,1)--(0,1)--cycle;
    \draw[vecblue!80!black,thick] (0,0)--(1,0)--(1,1)--(0,1)--cycle;
    \draw[->,vecblue,very thick] (0,0)--(1,0) node[below=1pt]{\scriptsize $\vec{e}_1$};
    \draw[->,vecred,very thick] (0,0)--(0,1) node[left=1pt]{\scriptsize $\vec{e}_2$};
    \filldraw[black] (0,0) circle (1.5pt);
    \node at (0.6,-1.1) {\small (a) Start: input space};
  \end{scope}

  % Arrow B
  \draw[->,gray!70,very thick] (1.9,0.7) -- (3.0,0.7) node[midway,above]{\small $\mat{B}$ (shear)};

  % Stage 2: After shear B = [1 1; 0 1]
  \begin{scope}[shift={(3.7,0)}]
    \draw[help lines,gray!25] (-0.3,-0.3) grid (2.4,1.6);
    \draw[->,gray!60] (-0.4,0)--(2.5,0);
    \draw[->,gray!60] (0,-0.4)--(0,1.7);
    \fill[quotientorange!20] (0,0)--(1,0)--(2,1)--(1,1)--cycle;
    \draw[quotientorange!80!black,thick] (0,0)--(1,0)--(2,1)--(1,1)--cycle;
    \draw[->,vecblue,very thick] (0,0)--(1,0) node[below=1pt]{\scriptsize $\mat{B}\vec{e}_1$};
    \draw[->,vecred,very thick] (0,0)--(1,1) node[above right=1pt]{\scriptsize $\mat{B}\vec{e}_2$};
    \filldraw[black] (0,0) circle (1.5pt);
    \node at (1.0,-1.1) {\small (b) Step 1: $\mat{B}$ acts};
  \end{scope}

  % Arrow A
  \draw[->,gray!70,very thick] (6.5,0.7) -- (7.8,0.7) node[midway,above]{\small $\mat{A}$ (rotate)};

  % Stage 3: After rotation A = [0 -1; 1 0]
  \begin{scope}[shift={(10.0,0)}]
    \draw[help lines,gray!25] (-1.4,-0.3) grid (0.6,2.3);
    \draw[->,gray!60] (-1.5,0)--(0.7,0);
    \draw[->,gray!60] (0,-0.4)--(0,2.4);
    \fill[structurepurple!20] (0,0)--(0,1)--(-1,2)--(-1,1)--cycle;
    \draw[structurepurple!80!black,thick] (0,0)--(0,1)--(-1,2)--(-1,1)--cycle;
    \draw[->,vecblue,very thick] (0,0)--(0,1) node[below right=1pt]{\scriptsize $\mat{A}\mat{B}\vec{e}_1$};
    \draw[->,vecred,very thick] (0,0)--(-1,1) node[above left=1pt]{\scriptsize $\mat{A}\mat{B}\vec{e}_2$};
    \filldraw[black] (0,0) circle (1.5pt);
    \node at (-0.3,-1.1) {\small (c) Step 2: $\mat{A}$ acts};
  \end{scope}

  % Composition curved arrow underneath
  \draw[->,structurepurple,very thick] (0.6,-1.6) to[out=-20,in=-160] node[midway,below=2pt]{\small Single composite transformation $\mat{A}\mat{B}$} (9.7,-1.6);
\end{tikzpicture}
\caption{Matrix multiplication represents the composition of linear transformations. Applying $\mat{B}$ (a shear)
followed by $\mat{A}$ (a $90^\circ$ rotation) transforms the space in two stages. The product matrix $\mat{A}\mat{B}$
encodes the transition from the original input grid (a) straight to the final rotated shape (c) in a single step.}
\label{ch01:composition}
\end{figure}

\begin{connectionbox}
This makes the set of invertible $n\times n$ matrices into a \emph{group} under multiplication -- the
general linear group $GL_n(\R)$. An invertible matrix is one whose transformation can be fully undone
without flattening space, and a group is a mathematical collection of actions that can be combined and
reversed. Here, the group operation is ``compose these two motions of space.'' Rotations alone
form a smaller group; rotations and reflections together form another. Geometry, in the program
launched by Felix Klein in 1872, \emph{is} the study of which properties survive a group
of transformations. Linear algebra is therefore the gateway to symmetry, and the symmetry groups it
houses run all the way through crystallography, particle physics, and computer graphics.
\end{connectionbox}

\section{Where This All Leads}

We have planted three seeds: vectors, the combinations you build from them (span), and the linear
transformations that move them around (matrices). Almost every later chapter is one of these seeds
grown tall. Here is the shape of what is coming.

\begin{itemize}
  \item \textbf{Solving systems and the four subspaces.} Elimination, rank, column space and null
  space -- the complete answer to ``how many solutions does $\mat A\vec x=\vec b$ have, and why.''
  Every question lives in the row/column duality of this chapter.
  \item \textbf{Orthogonality and projection.} Orthogonality is the rigorous geometric word for
  perpendicularity -- meeting at a clean $90^\circ$ right angle. Projection asks: what is the closest
  point in a subspace to a given vector? This single question is the engine of least-squares fitting
  in applied statistics.
  \item \textbf{Determinants.} One number that reads off the area- or volume-scaling factor of a
  transformation -- and tells you instantly whether it can be undone.
  \item \textbf{Eigenvalues and eigenvectors.} The special directions a transformation merely
  stretches without turning. They explain how systems evolve over time, why some structures resonate,
  and how Google originally ranked the web.
  \item \textbf{The Singular Value Decomposition.} A foundational theorem of matrix
  analysis: \emph{every} matrix factors into rotate, stretch, rotate (Figure~\ref{ch01:svd}).
  The SVD is the mathematical engine of data compression, recommendation systems, and principal
  component analysis.
\end{itemize}

\begin{figure}[H]
\centering
\begin{tikzpicture}[scale=0.88,>=Stealth]
  % stage 0: unit circle with two orthogonal axes
  \begin{scope}
    \draw[gray!55,->] (-1.3,0)--(1.3,0); \draw[gray!55,->] (0,-1.3)--(0,1.3);
    \draw[vecblue,very thick] (0,0) circle (1);
    \draw[->,vecred,thick] (0,0)--(0.71,0.71);
    \draw[->,structurepurple,thick] (0,0)--(-0.71,0.71);
    \node at (0,-1.8){\small unit circle};
  \end{scope}
  \draw[->,gray,thick] (1.7,0)--(2.5,0) node[midway,above]{\small $V\T$};
  % stage 1: rotated circle (still a circle)
  \begin{scope}[xshift=4.2cm]
    \draw[gray!55,->] (-1.3,0)--(1.3,0); \draw[gray!55,->] (0,-1.3)--(0,1.3);
    \draw[vecblue,very thick] (0,0) circle (1);
    \draw[->,vecred,thick] (0,0)--(1,0);
    \draw[->,structurepurple,thick] (0,0)--(0,1);
    \node at (0,-1.8){\small rotate};
  \end{scope}
  \draw[->,gray,thick] (5.9,0)--(6.7,0) node[midway,above]{\small $\Sigma$};
  % stage 2: stretch into ellipse along axes
  \begin{scope}[xshift=8.4cm]
    \draw[gray!55,->] (-1.9,0)--(1.9,0); \draw[gray!55,->] (0,-1.3)--(0,1.3);
    \draw[vecblue,very thick] (0,0) ellipse (1.7 and 0.7);
    \draw[->,vecred,thick] (0,0)--(1.7,0);
    \draw[->,structurepurple,thick] (0,0)--(0,0.7);
    \node at (0,-1.8){\small stretch};
  \end{scope}
  \draw[->,gray,thick] (10.5,0)--(11.3,0) node[midway,above]{\small $U$};
  % stage 3: final rotated ellipse
  \begin{scope}[xshift=13.4cm]
    \draw[gray!55,->] (-1.9,0)--(1.9,0); \draw[gray!55,->] (0,-1.9)--(0,1.9);
    \draw[vecblue,very thick,rotate=25] (0,0) ellipse (1.7 and 0.7);
    \draw[->,vecred,thick] (0,0)--({1.7*cos(25)},{1.7*sin(25)});
    \draw[->,structurepurple,thick] (0,0)--({-0.7*sin(25)},{0.7*cos(25)});
    \node at (0,-2.3){\small rotate again};
  \end{scope}
\end{tikzpicture}
\caption{A preview of the Singular Value Decomposition $\mat A = U\Sigma V\T$. Every matrix, however
complicated, does only this to space: rotate ($V\T$), stretch along axes ($\Sigma$), rotate again
($U$). The unit circle always lands on an ellipse, and its axis lengths are the singular values.}
\label{ch01:svd}
\end{figure}

\begin{applicationbox}
\textbf{Linear algebra in machine learning.} A feedforward neural network computes a sequence
of affine transformations interleaved with coordinate-wise nonlinearities. Each layer takes its input
vector $\vec x$, computes $\mat W\vec x + \vec b$ -- a linear transformation followed by a shift -- and
then applies a simple nonlinear function such as $\max(0,\cdot)$ to each entry.
The weight matrices $\mat W$ store the learned parameters; training adjusts the numerical entries of
these matrices via gradient descent. When such a model has billions of parameters, essentially all of
them live inside matrices, and evaluation requires multiplying matrices by vectors at scale. The
attention mechanism in transformer architectures is built from weight matrices ($\mat Q$, $\mat K$,
$\mat V$) and matrix products. GPUs power modern deep learning primarily because they are specialized
hardware for parallel matrix multiplication.
\end{applicationbox}

The same vocabulary reaches everywhere else too. In computer graphics, every rotation and zoom of a
3D scene is a matrix acting on vertices. In networks, the connections of a graph form a matrix whose
eigenvectors reveal communities and importance. In statistics, a covariance matrix's eigenvectors are
the principal directions of your data. In signal processing, the Fourier transform is a
change of basis into sinusoidal harmonics.

\begin{aestheticbox}
By the end of this book, you will be able to look at a system of equations, a high-dimensional
dataset, a geometric transformation, or an optimization step, and recognize the same mathematical
structure: vectors combined by weights, and linear maps transforming space.
\end{aestheticbox}

\begin{historybox}
Linear algebra grew out of a two-century search for the right algebra to describe geometry.
In 1637, René Descartes published \emph{La Géométrie}, introducing coordinate geometry: curves and lines
were converted into equations involving real numbers $(x,y,z)$. While Descartes' coordinates were a colossal
breakthrough, they came with a cost: geometry was reduced to mechanical algebra on coordinate numbers,
obscuring the intrinsic geometric objects -- the directed arrows, the planes, the physical forces -- under
a mountain of coordinate-dependent algebra.

Two centuries later, mathematicians searched for a direct algebra of space. On October 16, 1843, while walking
along the Royal Canal in Dublin, Sir William Rowan Hamilton had a flash of inspiration and carved his
equations for \emph{quaternions} into the stone of Brougham Bridge: $i^2 = j^2 = k^2 = ijk = -1$. Quaternions
are a four-dimensional number system that extends the complex numbers by adding three imaginary units to encode
three-dimensional rotations directly. But they were complicated hybrids, mixing a single real number (scalar part)
with a three-dimensional spatial direction (vector part) into an unwieldy package. In parallel, the German
schoolteacher Hermann Grassmann published his \emph{Ausdehnungslehre} (Theory of Extension) in 1844,
inventing vector spaces of arbitrary dimension and abstract geometric products, though his dense philosophical style was so
difficult that contemporaries ignored it for decades. In 1858, Arthur Cayley formalized the \textbf{matrix} as an
independent algebraic object, showing that multiplying rectangular arrays corresponds to composing linear transformations.

In the 1880s, the American physicist Josiah Willard Gibbs at Yale and the English electrical engineer Oliver
Heaviside independently realized that James Clerk Maxwell's theory of electromagnetism was weighed down by
quaternion notation. Gibbs and Heaviside split Hamilton's quaternion product into two separate operations: they discarded
the full quaternion product and extracted the \emph{dot product} (which multiplies two vectors to yield a single
number measuring how much they point together) and the \emph{cross product} (which takes two three-dimensional vectors
and produces a perpendicular vector whose length is the area of the parallelogram between them).

A ferocious controversy erupted, known as the \emph{Great Vector-Quaternion War} of the 1890s. Peter Guthrie Tait,
Hamilton's fiercest disciple, attacked Gibbs and Heaviside in the pages of \emph{Nature}, denouncing their vector
system as a ``monstrous hybrid'' and a ``mutilated quaternion.'' Heaviside fired back with characteristic wit:
``Quaternions were born with a great deal of promise, but they never made good... Vector analysis is what the
physical mathematician wants, in his daily work with electromagnetic fields, fluid dynamics, and mechanics.''
Heaviside and Gibbs prevailed: vector analysis became the universal language of physics and engineering.

Treating vectors as primary objects and matrices as linear transformations of space transformed
modern mathematics. Instead of drowning in separate scalar equations for every coordinate, an entire coupled
system across hundreds of variables can be written in four symbols: $\mat{A}\vec{x} = \vec{b}$. Linear
transformations provide an exact geometric framework for rotations, shears, projections, and deformations
of space. Yet relying solely on the mental picture of spatial arrows hits a wall the moment we venture beyond
three dimensions into image processing, econometrics, or neural networks with millions of parameters.
To harness the geometry of high-dimensional space without getting lost, we need formal algebraic coordinates
and systematic algorithms to solve coupled constraints. That bridge begins in Chapter~\ref{ch02:arrow},
which formalizes vectors in $\R^n$ by uniting the arrow, the coordinate list, and the data point, before
Chapter~\ref{ch03:def-lincomb} develops linear combinations and span.
\end{historybox}

\begin{focusbox}
\textbf{The three ideas to carry out of this chapter.}
\begin{enumerate}
  \item A \textbf{vector} is anything you can scale and add; arrows in the plane are the picture, but
  images, functions, and solution sets are vectors too.
  \item A \textbf{linear combination} mixes vectors by scaling and adding; the set of all reachable
  combinations is the \textbf{span}, and $\mat A\vec x=\vec b$ asks whether $\vec b$ is in the span of
  $\mat A$'s columns.
  \item A \textbf{matrix is a linear transformation} -- a verb that moves space, recorded by where it
  sends the basic directions. Its columns \emph{are} those destinations.
\end{enumerate}
Hold these three and the rest of the book is elaboration.
\end{focusbox}

Before turning the page to Chapter~2, test these three ideas on the exercises below.

\section{Exercises}

\begin{exercise}[Warmup: scale and add]
Let $\vec u = \colvec{1\\-2}$ and $\vec v = \colvec{3\\1}$. Compute $2\vec u + \vec v$,
$\vec u - 3\vec v$, and $\tfrac12(\vec u + \vec v)$. Sketch $\vec u$, $\vec v$, and $2\vec u + \vec v$
as arrows from the origin, and confirm the tip-to-tail picture for the sum.
\end{exercise}

\begin{exercise}[Spans in the plane]
For each pair, decide whether the span is a line or all of $\R^2$, and justify geometrically.
\[
\text{(a) } \colvec{1\\2},\ \colvec{2\\4};\qquad
\text{(b) } \colvec{1\\0},\ \colvec{0\\1};\qquad
\text{(c) } \colvec{1\\1},\ \colvec{-1\\1}.
\]
\end{exercise}

\begin{exercise}[Both pictures of one system]
Consider $\begin{aligned} x + 2y &= 4 \\ 3x - y &= 5 \end{aligned}$
\begin{enumerate}
  \item Draw the row picture (two lines) and estimate the intersection.
  \item Write the column picture: which combination of $\colvec{1\\3}$ and $\colvec{2\\-1}$ gives
  $\colvec{4\\5}$?
  \item Solve algebraically and check that the row and column answers agree.
\end{enumerate}
\end{exercise}

\begin{exercise}[Reading a matrix as a verb]
The matrix $\mat A = \begin{bmatrix} 0 & -1 \\ 1 & 0 \end{bmatrix}$ sends $\vec e_1=\colvec{1\\0}$ and
$\vec e_2=\colvec{0\\1}$ to which vectors? Apply $\mat A$ to $\colvec{2\\1}$ and to $\colvec{3\\0}$.
What familiar motion of the plane does $\mat A$ perform? Sketch the unit square and its image to
confirm.
\end{exercise}

\begin{exercise}[Identify the transformation]
For each matrix, describe in words what it does to the plane (stretch, rotate, reflect, shear,
collapse) by looking at where it sends $\vec e_1$ and $\vec e_2$:
\[
\text{(a) } \begin{bmatrix} 3 & 0 \\ 0 & 3\end{bmatrix},\quad
\text{(b) } \begin{bmatrix} 1 & 0 \\ 2 & 1\end{bmatrix},\quad
\text{(c) } \begin{bmatrix} 0 & 1 \\ 1 & 0\end{bmatrix},\quad
\text{(d) } \begin{bmatrix} 1 & 1 \\ 1 & 1\end{bmatrix}.
\]
One of these flattens the whole plane onto a line. Which, and onto which line?
\end{exercise}

\begin{exercise}[Functions are vectors]
Let $V$ be the set of all functions of the form $f(x) = a + b\cos x + c\sin x$.
\begin{enumerate}
  \item Show that if $f,g\in V$ then $f+g\in V$ and $2f\in V$, so $V$ is closed under scaling and
  adding.
  \item Argue that $V$ behaves like $\R^3$: each function corresponds to the triple $(a,b,c)$. Which
  three ``basic direction'' functions play the role of $\vec e_1,\vec e_2,\vec e_3$?
\end{enumerate}
\end{exercise}

\begin{exercise}[Linearity, tested]
Decide whether each map $T:\R^2\to\R^2$ is linear by checking the two promises $T(c\vec x)=cT(\vec x)$
and $T(\vec x+\vec y)=T(\vec x)+T(\vec y)$.
\[
\text{(a) } T\colvec{x\\y}=\colvec{2x\\x-y};\quad
\text{(b) } T\colvec{x\\y}=\colvec{x+1\\y};\quad
\text{(c) } T\colvec{x\\y}=\colvec{xy\\0}.
\]
For each linear one, write down its matrix. For each nonlinear one, name the exact promise it breaks.
\end{exercise}

\begin{exercise}[Composition by hand]
Let $\mat R=\begin{bmatrix}0&-1\\1&0\end{bmatrix}$ (rotate $90^\circ$) and
$\mat F=\begin{bmatrix}1&0\\0&-1\end{bmatrix}$ (flip vertically). Apply ``$\mat F$ then $\mat R$'' to
$\colvec{1\\0}$ and to $\colvec{0\\1}$ by tracking the vectors through both steps. Then apply
``$\mat R$ then $\mat F$.'' Do the two orders give the same result? What does your answer tell you
about whether matrix multiplication can be commutative?
\end{exercise}

\begin{exercise}[Proof: linearity forces the origin to stay put]
Suppose $T$ is linear. Using only the scaling promise, prove that $T(\vec 0) = \vec 0$.
(\emph{Hint:} $\vec 0 = 0\cdot\vec 0$.) Explain in one sentence why this means the map
$\vec x \mapsto \vec x + \colvec{1\\1}$ -- a pure shift -- can never be linear.
\end{exercise}

\begin{exercise}[Challenge: when do three vectors fill space?]
You hold three vectors $\vec a,\vec b,\vec c$ in $\R^3$. Argue geometrically that their span is all of
$\R^3$ exactly when none of them lies in the plane spanned by the other two. Give an explicit example
of three vectors whose span is only a plane, and verify that one of them is a linear combination of the
other two. (You are rediscovering linear dependence -- the subject of a coming chapter.)
\end{exercise}

\begin{exercise}[Challenge: a recipe is unique when directions differ]
Suppose $\vec v$ and $\vec w$ in $\R^2$ point in different directions (neither is a multiple of the
other). Prove that any target $\vec b$ can be written as $c\vec v + d\vec w$ in \emph{exactly one}
way. (\emph{Hint:} if $c\vec v + d\vec w = c'\vec v + d'\vec w$, move everything to one side and use
that $\vec v,\vec w$ are not parallel.) Connect this to why the column picture of a $2\times2$ system
with independent columns has a unique solution.
\end{exercise}
